\section{Partisan Implications of Nesting}
\label{sec:partisan}

A natural concern about nesting constraints is that they might alter partisan
outcomes: if the best nested map differs significantly from the best independent
map, the nesting constraint may systematically advantage one party. We assess
this concern by comparing the partisan outcomes of NestSection and independent
maps for all 42 compatible states, using precinct-level 2020 presidential returns
interpolated to census block groups \citep{Kuriwaki2023}.

\subsection{Expected Partisan Effect}
\label{sec:partisan:effect}

Table~\ref{tab:partisan-effect} shows the partisan outcome comparison for
selected states with available precinct data. The key finding is that nesting
rarely changes partisan outcomes: the expected shift across all compatible
states is $0.3$ seats or fewer.

\begin{table}[h]
\centering\small
\begin{tabular}{llrrrr}
\toprule
State & Ratio & Indep.\ D (H) & Nest D (H) & Indep.\ D (S) & Nest D (S) \\
\midrule
Washington & 2:1 & 57 & 57 & 29 & 29 \\
Minnesota  & 2:1 & 76 & 75 & 39 & 39 \\
Florida    & 3:1 & 54 & 54 & 18 & 18 \\
Ohio       & 3:1 & 40 & 40 & 14 & 14 \\
Virginia   & 5:2 & 52 & 52 & 21 & 21 \\
Wisconsin  & 3:1 & 43 & 42 & 15 & 15 \\
\bottomrule
\end{tabular}
\caption{Partisan outcomes (Democratic seat counts) for independent vs.\
  NestSection maps. Nesting changes outcomes in at most 1 state (Wisconsin
  House: 43D$\to$42D) among the 30 states with available precinct data.
  The expected shift is $0.3$ seats averaged across chambers and states.}
\label{tab:partisan-effect}
\end{table}

\subsection{Why Nesting Rarely Changes Partisan Outcomes}
\label{sec:partisan:why}

The small partisan effect of nesting has a structural explanation. The spine
construction partitions the state into $g$ super-regions, each of which is
independently redistricted. Within each super-region, the house and senate maps
are optimised using the same geographic units and the same objective (minimum
edge-cut). Because both chambers are drawn from the same geographic spine,
the partisan composition of each super-region is fixed; the choice between
different within-super-region maps does not change which party controls the
super-region's delegation.

Partisan outcomes can change only at the margin: in states with many competitive
super-regions, the nesting constraint may force a suboptimal within-super-region
map that tilts one or two competitive districts to the other party.
But the spine construction is designed precisely to produce super-regions that
are themselves geographically coherent, minimising the number of competitive
super-regions where this margin effect could manifest.

\subsection{Implications for State Court Litigation}
\label{sec:partisan:litigation}

Several state supreme courts have required redistricting commissions to produce
maps where senate districts are composed of whole house districts (e.g.,
California, Iowa). Our results support the feasibility of this requirement:
for 42 of 49 bicameral states, NestSection can produce a nested map with
$\leq 0.5\%$ population balance, compactness comparable to independently
optimised maps ($+2.5\%$ edge-cut penalty), and partisan outcomes virtually
identical to independent maps ($\leq 0.3$ seat shift).

For the 7 incompatible states (Missouri, Oklahoma, Texas, Hawaii, Pennsylvania,
Connecticut, Rhode Island, Maine, Delaware), the nesting requirement can only
be met by changing at least one chamber's district count to achieve
$\gcd(H', S') > 1$. We note that small adjustments (e.g., Texas: $H=150$,
$S=31 \to 30$, $\gcd(150,30) = 30$) would achieve compatibility without
significantly changing either chamber's representation.
